\section{Results}
\label{sec:results}

\subsection{Offline: joint decoding against cascades}
\label{sec:res-offline}

On the AMI reference windows (Table~\ref{tab:main7full}), \sys{} (offline) reaches
\num{v7:tx:en2zh:O2:cp} / \num{v7:tx:zh2en:O2:cp} cpBLEU$_\text{sess}$ (En$\rightarrow$Zh / Zh$\rightarrow$En) against
\num{v7:tx:en2zh:C-ft2:cp} / \num{v7:tx:zh2en:C-ft2:cp} for the cascade whose translator was trained on our windows
(exploratory comparison) and \num{v7:tx:en2zh:C-turn:cp} / \num{v7:tx:zh2en:C-turn:cp} for the one retrained on its own
diarizer turns (confirmatory: \sys{} ahead by \num{v7:reg:E-cas1:diff} and \num{v7:reg:E-cas2:diff}, $p_\text{Holm}$
\num{v7:reg:E-cas1:holm}). A speaker-blind model relabelled after decoding with the same diarizer reaches
\num{v6:relabel:en2zh:cp} / \num{v6:relabel:zh2en:cp}, close to \sys{} overall.

\subsection{Overlapped speech}
\label{sec:res-overlap}

\begin{table*}[t]
\centering\small
\setlength{\tabcolsep}{4pt}
\begin{tabular}{lrrrrrrrrr}
\toprule
 & \multicolumn{3}{c}{En$\rightarrow$Zh} & \multicolumn{3}{c}{Zh$\rightarrow$En} & \multicolumn{3}{c}{AliMeeting} \\
\cmidrule(lr){2-4}\cmidrule(lr){5-7}\cmidrule(lr){8-10}
System & all & overlap & CR & all & overlap & CR & all & overlap & CR \\
\midrule
\multicolumn{10}{l}{\emph{Reference windows, offline}} \\
Cascade, same data & 36.13 & 26.21 & 89.3 & 21.96 & 14.36 & 87.4 & 15.70 & 9.79 & 93.3 \\
Speaker-blind + relabel & 41.39 & 30.98 & 51.8 & 24.36 & 16.90 & 50.6 & 16.27 & 12.04 & 58.0 \\
\sys{} (offline) & 42.06 & 34.56 & 76.3 & 25.43 & 19.72 & 79.0 & 15.19 & 12.10 & 79.2 \\
\multicolumn{10}{l}{\emph{Whole meetings, online, 1\,s reads}} \\
Offline model + LA2 & 34.30 & 22.34 & 73.4 & 23.94 & 18.88 & 78.7 & 11.11 & 7.19 & 75.6 \\
Simul model + LA2 & 39.31 & 32.91 & 75.5 & 23.96 & 17.89 & 79.4 & 11.91 & 7.60 & 76.7 \\
Speaker-blind Simul + causal relabel & 35.22 & 27.17 & 65.5 & 19.33 & 9.06 & 54.4 & 12.67 & 8.55 & 63.0 \\
\sys{}-Simul & 40.09 & 35.71 & 82.3 & 24.53 & 17.73 & 80.2 & 11.47 & 9.18 & 80.2 \\
\bottomrule
\end{tabular}
\caption{Overlap-aware evaluation (\S\ref{sec:eval-overlap}): tcpBLEU over all units and over overlapped units, and
concurrent-speaker recall (CR, \%). Relabelling attaches one diarizer speaker to each translated segment. All rows and
metrics: Table~\ref{tab:overlap8}. AliMeeting references: machine translations of the human transcripts (Appendix~\ref{app:qe7}).}
\label{tab:ovmain8}
\end{table*}


Table~\ref{tab:ovmain8} splits the scores by overlap (\S\ref{sec:eval-overlap}). On single-speaker units relabelling
matches \sys{} (\num{v8:ov:en2zh:O2-nosd-relabel:1spk} vs.\ \num{v8:ov:en2zh:O2:1spk} in En$\rightarrow$Zh). On
overlapped units \sys{} leads by \num{v8:x:ov-O2-rel-en2zh:diff} (En$\rightarrow$Zh, 95\% interval
\num{v8:x:ov-O2-rel-en2zh:ci}) and \num{v8:x:ov-O2-rel-zh2en:diff} (Zh$\rightarrow$En, \num{v8:x:ov-O2-rel-zh2en:ci}),
both exploratory, and ties on AliMeeting (\num{v8:x:ov-O2-rel-ali:diff}); the gap opens with the amount of overlap
(Figure~\ref{fig:overlap8}). Relabelling gives one speaker per segment, so when two people talk at once one of them
disappears or is merged into the other (concurrent-speaker recall \num{v8:ov:en2zh:O2-nosd-relabel:cr}\%, collapse
\num{v8:ov:en2zh:O2-nosd-relabel:coll}\%); \sys{}, reading both speakers' activity, writes both (\num{v8:ov:en2zh:O2:cr}\%,
\num{v8:ov:en2zh:O2:coll}\%; recall gain \num{v8:x:cr-O2-rel-en2zh:diff} / \num{v8:x:cr-O2-rel-zh2en:diff} /
\num{v8:x:cr-O2-rel-ali:diff}, all $p<0.001$). Relabelling \sys{}'s own text shows that this gain is mainly
attribution, not better translation of overlapped speech (Appendix~\ref{app:v8}). The cascade scores far lower on overlapped units but recovers more
concurrent pairs (\num{v8:ov:en2zh:C-ft2:cr}\%): \sys{} drops short interjections (Appendix~\ref{app:v8}).
\label{sec:res-real}
\label{sec:res-speech}

\subsection{When the translation is heard}
\label{sec:res-pa}
\label{sec:res-v10}

\begin{table*}[t]
\centering\small
\setlength{\tabcolsep}{3.5pt}
\resizebox{\textwidth}{!}{\begin{tabular}{llrrrrrrrrr}
\toprule
 & & \multicolumn{3}{c}{text tcpBLEU $\uparrow$} & speech & \multicolumn{5}{c}{latency (s) $\downarrow$} \\
\cmidrule(lr){3-5}\cmidrule(lr){7-11}
Test set & System & all & ov & 1spk & ASR-BLEU $\uparrow$ & text LAAL & speech start & LAAL$_\text{block}$ & LAAL$_\text{spread}$ & EndOffset \\
\midrule
\multirow{3}{*}{AMI Zh$\rightarrow$En} & Deadline baseline (B) & 25.02 & 19.97 & 27.10 & 18.07 & 2.99 & 2.99 & 3.12 & 4.01 & 4.08 \\
 & \sys{}: B + playback alignment & 25.02 & 19.97 & 27.10 & 20.35$^{*}$ & 2.99 & 2.96 & 3.03 & 3.61 & 3.13 \\
 & \sys{}+WA, $\bias=-1$ (\S\ref{sec:write}) & 25.34$^{\circ}$ & 20.53 & 27.30 & 20.45$^{*}$ & 3.07 & 3.04 & 3.10 & 3.69 & 3.20 \\
\midrule
\multirow{3}{*}{AliMeeting Zh$\rightarrow$En} & Deadline baseline (B) & 21.43 & 16.91 & 23.05 & 16.73 & 2.82 & 3.05 & 3.28 & 4.28 & 4.57 \\
 & \sys{}: B + playback alignment & 21.43 & 16.91 & 23.05 & 18.08$^{*}$ & 2.82 & 2.92 & 2.95 & 3.65 & 3.23 \\
 & \sys{}+WA, $\bias=-1$ (\S\ref{sec:write}) & 21.07$^{\circ}$ & 17.38 & 22.68 & 17.84$^{*}$ & 3.06 & 3.20 & 3.19 & 3.91 & 3.45 \\
\midrule
\multirow{3}{*}{AMI En$\rightarrow$Zh} & Deadline baseline (B) & 41.15 & 36.16 & 42.51 & 31.65 & 2.37 & 2.43 & 2.54 & 3.21 & 3.48 \\
 & \sys{}: B + playback alignment & 41.15 & 36.16 & 42.51 & 32.23$^{*}$ & 2.37 & 2.39 & 2.43 & 2.95 & 2.91 \\
 & \sys{}+WA, $\bias=-1$ (\S\ref{sec:write}) & 41.74$^{\circ}$ & 36.48 & 43.24 & 32.34$^{\circ}$ & 2.47 & 2.50 & 2.52 & 3.05 & 2.99 \\
\bottomrule
\end{tabular}}
\caption{Whole test meetings, run online, each system decoded once (\S\ref{sec:setup}); B and \sys{} share their text; the first test decoding of \sys{}+WA (before $\bias$) is in Table~\ref{tab:delta0}. tcpBLEU: all / overlapped / single-speaker units; ASR-BLEU: mean of two TTS seeds; latency: per speaker stream, chunk protocol, without computation. Against B (Table~\ref{tab:sig10}): $^{*}$ $p_\text{Holm}<0.05$, $^{\circ}$ not significant; latency intervals: Table~\ref{tab:postreview2}.}
\label{tab:main10}
\end{table*}


Table~\ref{tab:main10} gives the streaming systems on whole test meetings; all share training data, online windows
and diarizer, and B and \sys{} share their text. Playback alignment raises ASR-BLEU by
\num{v10:sig:test:ami:speech:e2-b:diff} / \num{v10:sig:test:en:speech:e2-b:diff} /
\num{v10:sig:test:ali:speech:e2-b:diff} (AMI Zh$\rightarrow$En / En$\rightarrow$Zh / AliMeeting; all
$p_\text{Holm}<0.001$, Table~\ref{tab:sig10}). It also changes \emph{when} speech is heard, on every latency definition (meeting-bootstrap intervals,
Appendix~\ref{app:postreview}): conventional block-start LAAL changes by
\num{v10:pr:test:ami:e2-b:bslaal:diff} \num{v10:pr:test:ami:e2-b:bslaal:ci} / \num{v10:pr:test:en:e2-b:bslaal:diff}
\num{v10:pr:test:en:e2-b:bslaal:ci} / \num{v10:pr:test:ali:e2-b:bslaal:diff} \num{v10:pr:test:ali:e2-b:bslaal:ci}\,s
and EndOffset by \num{v10:pr:test:ami:e2-b:endoff:diff} / \num{v10:pr:test:en:e2-b:endoff:diff} /
\num{v10:pr:test:ali:e2-b:endoff:diff}\,s, while speech start moves by only \num{v10:pr:test:ami:e2-b:startoff:diff} /
\num{v10:pr:test:en:e2-b:startoff:diff} / \num{v10:pr:test:ali:e2-b:startoff:diff}\,s. Speech ends earlier because
the PA Talker speaks faster (\num{v10:pr:test:ami:e2:rate} instead of \num{v10:pr:test:ami:b:rate} words per second
in AMI Zh$\rightarrow$En), so the same text takes \num{v10:pr:test:ami:e2-b:speechpct} /
\num{v10:pr:test:en:e2-b:speechpct} / \num{v10:pr:test:ali:e2-b:speechpct}\% less time and blocks queue less
(Table~\ref{tab:decomp10}). UTMOS rises slightly (\num{v10:pr:test:ami:b:utmos} $\rightarrow$
\num{v10:pr:test:ami:e2:utmos}), but pauses of at least 0.5\,s per minute of speech rise from
\num{v10:pr:test:ami:b:pauses} to \num{v10:pr:test:ami:e2:pauses} (AMI Zh$\rightarrow$En; similar elsewhere), a fluency cost
UTMOS does not capture\ifFluencyPA{}; the rate budget (\S\ref{sec:res-fluency}) addresses it\fi.
\ifMetricsV
As selection uses a CTC recogniser and ASR-BLEU uses ASR, the gain could partly reflect ASR-friendly speech; with a
second recogniser the gain is \num{v11:asr2:gain:ami} / \num{v11:asr2:gain:en} / \num{v11:asr2:gain:ali}
(Appendix~\ref{app:metricsv11}).
\else
As selection uses a CTC recogniser and ASR-BLEU uses ASR, the gain may partly reflect ASR-friendly speech.
\fi

\ifFactorial
\paragraph{Which part does it?}
\begin{table*}[t]
\centering\scriptsize
\setlength{\tabcolsep}{2pt}
\begin{tabular}{lrrrrrrrr}
\toprule
 & \multicolumn{3}{c}{ASR-BLEU} & \multicolumn{3}{c}{LAAL$_\text{block}$ (s)} & rate & pause \\
\cmidrule(lr){2-4}\cmidrule(lr){5-7}
Speech layer on B's text & Z$\rightarrow$E & E$\rightarrow$Z & Ali & Z$\rightarrow$E & E$\rightarrow$Z & Ali & w/s & /min \\
\midrule
R0 B & 18.07 & 31.64 & 16.73 & 3.12 & 2.54 & 3.28 & 2.40 & 3.3 \\
R1 + read-level merging & 18.32 & 31.66 & 16.61 & 3.09 & 2.50 & 3.16 & 2.46 & 3.5 \\
R2 + backlog rate rule & 19.88 & 31.91 & 17.48 & 3.04 & 2.45 & 3.00 & 3.20 & 6.8 \\
R3 + best-of-2 CTC & 20.27 & 32.26 & 17.81 & 3.04 & 2.45 & 2.99 & 3.21 & 6.9 \\
R3$^\prime$ R3, rate-matched & 20.24 & 32.17 & 17.85 & 3.03 & 2.43 & 2.94 & 3.38 & 7.7 \\
R4 R2 + PA Talker & 20.17 & 31.93 & 17.82 & 3.03 & 2.43 & 2.95 & 3.48 & 8.2 \\
R5 \sys{} (R4 + best-of-2) & 20.35 & 32.23 & 18.08 & 3.03 & 2.43 & 2.95 & 3.49 & 8.3 \\
\bottomrule
\end{tabular}
\caption{Factorial ablation of the playback layer on the test sets, B's text throughout, each row decoded once with two TTS seeds. R3$^\prime$ raises the rule's rate shift until its mean development rate matches \sys{}'s (shift 8 / 8; AMI Zh$\rightarrow$En rate 3.38 vs.\ 3.49). Rate and pauses: AMI Zh$\rightarrow$En.}
\label{tab:factorial}
\end{table*}

B+PA differs from B in merging, a backlog-dependent rate, best-of-two CTC selection and the learned Talker.
Table~\ref{tab:factorial} switches them on one at a time on B's text, with a rule-based layer whose mean rate is set to
the PA Talker's on development meetings. \num{v11:fact:verdict}
\else
B+PA differs from B in four ways (merging, the fine-tuned Talker, a backlog-dependent rate, best-of-two CTC
selection); on a fixed text a rule-based layer with the same rate and selection gets most of the gain, the learned layer
adding \num{v10:tk:ami:c-a:asr} / \num{v10:tk:ali:c-a:asr} ASR-BLEU (Table~\ref{tab:talker10}).
\fi

\ifFluencyPA
\paragraph{A rate budget.}
\label{sec:res-fluency}
\num{v11:flu:verdict}
\fi

\ifAliEval
\paragraph{Unused meetings.}
\label{sec:res-alieval}
To test the speech-layer claims without the selection effects of a test set used throughout development, we fixed the
systems and hypotheses in an analysis plan and decoded once the eight AliMeeting evaluation meetings, never used before,
scored against the human English translations released with DiariST (Table~\ref{tab:alieval}, Appendix~\ref{app:alieval}).
\num{v11:ae:verdict}
\fi

\ifWallclock
\paragraph{Computation.}
\label{sec:res-compute}
\num{v11:wc:verdict} (Appendix~\ref{app:wallclock}).
\fi

\subsection{Test-set hygiene}
\label{sec:res-hygiene}

\sys{} and B are a fixed-step Thinker with Talkers chosen on development data; no choice was made on test. Still, the
same 16 AMI and 20 AliMeeting test meetings served throughout development while windows, latency protocol, merging and
the AliMeeting references changed, so adaptive overfitting cannot be ruled out\ifAliEval{} (the unused meetings above
address this for the speech layer)\fi; the write-alignment rows carry two further test-informed decisions
(\S\ref{sec:write}). Analysis plans were internal (Table~\ref{tab:reg9}).
